\section{Positive-Only Separator Add-On}\label{sec:separator}

Set
\[
  X_n=\BtildeI_n(1,1).
\]
This section uses the equivalent one-indexed signed-word convention
$w=(w_1,\ldots,w_n)$, with $|w_1|,\ldots,|w_n|$ a permutation of $\{1,\ldots,n\}$ and
$\rho(i)=n+1-i$.  Thus
\[
  X_n=\{w\in B_n:\#\{i:w_i=i\}=1,\ \#\{i:w_i=\rho(i)\}=1\}.
\]
For a signed word define the Reiner-style total order \cite{Reiner1993}
\[
  T(w): -w_n<\cdots<-w_1<0<w_1<\cdots<w_n.
\]
Let $\Sigma_n=\{-n,\ldots,-1,0,1,\ldots,n\}$.  For $X\subseteq B_n$, put
\[
  C(X)=\bigcap_{w\in X}\{(a,b)\in\Sigma_n^2:a\text{ precedes }b\text{ in }T(w)\}.
\]
If $X=L_B(P)$ is the signed linear-extension set of a signed poset $P$, then every strict relation of $P$, and every
relation in its transitive closure, lies in $C(X)$.

\begin{theorem}[Positive-only separator]\label{thm:s12}
For every $n\ge4$, the layer $X_n=\BtildeI_n(1,1)$ is not a Reiner-style signed-poset linear-extension set and is not
an exact left-translate tiler of $B_n$.  For $n=3$, the divisibility obstruction disappears and there is an exact
left-translate tiling of $B_3$ by $X_3$.
\end{theorem}

\begin{center}
\begingroup
\footnotesize
\begin{tabular}{>{\raggedright\arraybackslash}p{0.28\textwidth}>{\raggedright\arraybackslash}p{0.39\textwidth}>{\raggedright\arraybackslash}p{0.22\textwidth}}
\textbf{Part} & \textbf{Proof type} & \textbf{Replay route}\\
\hline
signed-poset obstruction & local sign-extension construction plus exact finite $n=4$ bridge & S12A\\
$n=3$ tiling exception & exact finite cover & S12B\\
translate-tiler obstruction & rook-polynomial count, divisibility, tail inequalities, and finite window & S12B\\
\end{tabular}
\endgroup
\end{center}

\noindent\textit{Proof-route map.}  This is an additional scoped theorem, separate from the two-standard-channel rank theorem.
\begin{proof}
We first prove the signed-poset obstruction.  For $n\ge5$, any prescribed signs on at most two absolute labels can be
extended to an element of $X_n$.  Indeed, positive labels create fixed or reversal hits only when placed in their two
dangerous rows, while negative labels create no positive hit.  Place the prescribed positive labels, then add exactly
one positive fixed hit and one positive reversal hit; all remaining labels are placed negatively.  If the odd center is
used, it is used as the intended center hit, since the center supplies both hit types at once.  Otherwise auxiliary
positive labels are chosen away from the center and away from their dangerous rows.  With at most two prescribed labels,
there is always such a spare choice for $n\ge5$.  The odd center is the only place where $\rho(i)=i$ can interfere with the sign-extension construction; the construction either uses it deliberately as the single center hit or avoids it when choosing the auxiliary positive labels.

Given any ordered pair $a\ne b$ in $\Sigma_n$, choose signs so that $b$ lies before $0$ and $a$ lies after $0$ in
$T(w)$.  If $a,b$ are nonzero and $a\ne-b$, prescribe $|a|$ with the sign of $a$ and $|b|$ with the opposite sign to
$b$.  If $a=-b$, prescribe only $|a|$ with the sign of $a$.  If one entry is $0$, prescribe the other label on the
required side of $0$.  The sign-extension construction then gives $w\in X_n$ in which $b$ precedes $a$.  Hence
$C(X_n)=\varnothing$ for $n\ge5$.

The boundary case $n=4$ is finite: an exact finite check records witnesses in $X_4$ reversing all $72$ ordered pairs
in $\Sigma_4$.  Thus $C(X_4)=\varnothing$ as well.  If $X_n=L_B(P)$ for some signed poset $P$, all relations of $P$
would lie in $C(X_n)$, so $P$ would be the antichain and $L_B(P)=B_n$.  But $X_n\ne B_n$, for instance the all-negative
identity has no positive fixed or reversal hit.  This proves the signed-poset obstruction for all $n\ge4$.

For the tiler obstruction, an exact left-translate tiling by $X_n$ requires $N_n=|X_n|$ to divide
$G_n=|B_n|=2^n n!$.  The count of $X_n$ is obtained from the rook polynomial
\[
  R_{2m}(A,B,Z)=(1+2AZ+2BZ+A^2Z^2+B^2Z^2)^m,
\]
\[
  R_{2m+1}(A,B,Z)=R_{2m}(A,B,Z)(1+AZ+BZ+ABZ),
\]
where $A$ marks selected positive fixed events, $B$ selected positive reversal events, and $Z$ distinct forced cells.
Inclusion-exclusion for exactly one event of each type gives
\[
  N_n=\sum_{p,q\ge1,u}(-1)^{p+q}pq[A^pB^qZ^u]R_n(A,B,Z)2^{n-u}(n-u)!.
\]
Equivalently, for
\[
  S_n(Z)=\left.(A\partial_A)(B\partial_B)R_n(A,B,Z)\right|_{A=B=-1},
\]
we have
\[
  \frac{N_n}{G_n}=\sum_u \frac{[Z^u]S_n(Z)}{2^u(n)_u}.
\]
The differentiated series are
\[
S_{2m}(Z)=4m(m-1)Z^2(1-Z)^2(1-4Z+2Z^2)^{m-2},
\]
\[
S_{2m+1}(Z)=Z(1-4Z+2Z^2)^m+4m(m-1)Z^2(1-Z)^3(1-4Z+2Z^2)^{m-2}.
\]
Substituting $Z=-x$ makes the tails alternating.  The certificate
\path{P15_S12_GATE_B_TILER_CERTIFICATE}, routed in Appendix~\ref{app:certs}, records the exact finite window
$4\le n\le21$, the $n=3$ exact cover, and the symbolic all-$n$ tail inequalities below.

For even $n=2m\ge22$, write
\[
  S_{2m}(-x)=4m(m-1)x^2(1+x)^2f(x)^{m-2},\qquad f(x)=1+4x+2x^2.
\]
Because $f(x)=(1+(2+\sqrt2)x)(1+(2-\sqrt2)x)$ has positive linear factors, the coefficients of
$(1+x)^2f(x)^{m-2}$ are ultra-log-concave by Newton's inequalities.  If $c_u=[x^u]S_{2m}(-x)$, this gives
\[
  \frac{c_{u+1}}{2^{u+1}(2m)_{u+1}}\le
  \frac{c_u}{2^u(2m)_u},
\]
so the normalized alternating tail is decreasing.  The certificate then checks the two finite symbolic inequalities
$L_9>1/11$ and $U_6<1/10$ after the shift $m=t+11$; all coefficients of the resulting polynomials are positive.

For odd $n=2m+1\ge23$, the first coefficient of $S_{2m+1}(-x)$ is negative and all later coefficients are positive.
Indeed
\[
S_{2m+1}(-x)=x f(x)^{m-2}Q_m(x),
\]
and the certificate lists the coefficients of $Q_m$: the nonconstant coefficients are positive for $m\ge11$, while the
top three support degrees are supplied by the positive high coefficients of $Q_m$.  To prove the normalized tail is
monotone from degree $2$, put $C_m(x)=S_{2m+1}(-x)$ and use the equivalent coefficient positivity of
\[
  H_m(x)=2((2m+1)C_m(x)-xC_m'(x))-C_m(x)/x.
\]
The identity $H_m(x)=f(x)^{m-3}P_m(x)$ is checked symbolically, and Newton log-concavity for $f(x)^{m-3}$ reduces the
possible negative contribution of the linear term of $P_m$ to a positive quadratic lower bound for every $m\ge11$.
As in the even case, the certificate checks $L_9>1/11$ and $U_6<1/10$ after the shift $m=t+11$ by coefficient
positivity.  Thus the alternating-tail bounds and exact finite checks give
\[
  9< G_4/N_4<10,
\]
\[
  10<G_n/N_n<11\quad(n\ge5\text{ odd}),
\]
\[
  11<G_n/N_n<12\quad(6\le n\le20\text{ even}),
\]
\[
  10<G_n/N_n<11\quad(n\ge22\text{ even}).
\]
Thus $G_n/N_n$ is never an integer for $n\ge4$, so $X_n$ cannot tile $B_n$ by exact left translates.  For $n=3$, an exact finite check records an exact cover by $24$ left translates of the two-element set $X_3$.
\end{proof}
